% Preámbulo de los parciales · XeLaTeX · mismo estilo de los labs (rl_metodos/docs/curso/tex/preambulo.tex)
% Tipografía: Inter (cubre subíndices y griegas Unicode), con DejaVu Sans para los pocos símbolos que le faltan.
\documentclass[11pt,a4paper]{article}
\usepackage[margin=20mm,top=18mm,bottom=20mm]{geometry}
\usepackage{fontspec}
\setmainfont{Inter}[BoldFont={Inter SemiBold}, ItalicFont={Inter Italic}, BoldItalicFont={Inter SemiBold Italic}, Scale=0.95]
\setsansfont{Inter}
\setmonofont{DejaVu Sans Mono}[Scale=0.82]
\newfontfamily\dv{DejaVu Sans}[Scale=0.9]
\usepackage{unicode-math}
\setmathfont{latinmodern-math.otf}
\usepackage{polyglossia}\setdefaultlanguage{spanish}
\usepackage{xcolor,booktabs,longtable,array,calc,tabularx,enumitem,parskip,microtype,titlesec,fancyhdr,newunicodechar}
\usepackage[most]{tcolorbox}
\usepackage[hidelinks]{hyperref}
\emergencystretch=3em
\providecommand{\tightlist}{\setlength{\itemsep}{1pt}\setlength{\parskip}{0pt}}
\setcounter{secnumdepth}{-1}

\newunicodechar{∈}{{\dv ∈}}
\newunicodechar{∎}{{\dv ∎}}
\newunicodechar{⟨}{{\dv ⟨}}
\newunicodechar{⟩}{{\dv ⟩}}
\newunicodechar{☐}{{\dv ☐}}
\newunicodechar{✔}{{\dv ✔}}
\newunicodechar{✘}{{\dv ✘}}
\AtBeginDocument{\def\square{\mbox{\dv ☐}}}

\definecolor{tinta}{HTML}{1F2933}
\definecolor{acento}{HTML}{B7410E}
\definecolor{suave}{HTML}{F4F1EC}
\definecolor{gris}{HTML}{6B7280}
\definecolor{linea}{HTML}{D6D0C6}
\color{tinta}

\titleformat{\section}{\Large\bfseries\color{acento}}{}{0em}{}
\titleformat{\subsection}{\large\bfseries}{}{0em}{}
\titleformat{\subsubsection}{\normalsize\bfseries}{}{0em}{}
\titlespacing*{\section}{0pt}{1.4em}{0.4em}
\titlespacing*{\subsection}{0pt}{1.0em}{0.25em}
\setlist{nosep, leftmargin=1.4em, itemsep=2pt}
\renewcommand{\arraystretch}{1.3}
\setlength{\LTpre}{6pt}\setlength{\LTpost}{6pt}

% citas: el escenario inventado y las notas
\renewenvironment{quote}{\begin{tcolorbox}[colback=suave, colframe=linea, boxrule=0.5pt, arc=2pt, left=6pt, right=6pt, top=4pt, bottom=4pt]}{\end{tcolorbox}}
% código corto
\let\oldverbatim\verbatim
\let\endoldverbatim\endverbatim

% espacios para escribir
\newcommand{\hueco}[1]{\rule[-1.5pt]{\dimexpr #1\dimexpr 0.33em\relax\relax}{0.45pt}}
\newcommand{\lineasrespuesta}[1]{\par\vspace{2pt}{\color{linea}\foreach\i in {1,...,#1}{\vspace{0.75cm}\hrule height 0.5pt}}\par\vspace{4pt}}
\newcommand{\filarespuesta}{\rule{0pt}{1.25cm}}

\pagestyle{fancy}\fancyhf{}
\renewcommand{\headrulewidth}{0pt}
\fancyfoot[C]{\small\color{gris}\docpie\ · página \thepage}

\newcommand{\cabecera}[3]{%
  {\huge\bfseries\color{acento}#1\par}\vspace{4pt}
  {\large #2\par}\vspace{2pt}
  {\small\color{gris}#3\par}\vspace{6pt}
  {\color{linea}\hrule height 0.6pt}\vspace{8pt}}
